\section{Validation}\label{sec:mechanization}

The Lean development, the programs described in this section and their outputs are available at
\url{https://github.com/audenvenn/layouts-and-carries}.

\subsection{Lean formalisation}\label{sec:lean}

The core of \cref{sec:carry,sec:normalform,sec:composition} is formalised in Lean~4~\cite{moura2021}
with Mathlib~\cite{mathlib2020} (Lean and Mathlib \code{v4.37.0-rc1}); the development is about
$1\,500$ lines, contains no \code{sorry}, and depends only on the standard axioms \code{propext},
\code{Classical.choice} and \code{Quot.sound}. A flat layout is a list of pairs (size, stride)
with strides in an arbitrary additive commutative group, \code{eval} is the extended layout
function defined by recursion on the list exactly as in \cref{def:layout}, and a ``floor sum'' is
the function $x\mapsto\sum_j\floor{x/Q_j}\cdot\jump_j$ attached to a list of (scale, jump) pairs.
\Cref{tab:lean} lists what is proved. Statements marked \mech{} in the text are those of the
table; mode sizes are assumed positive throughout, as in \cref{def:shape}.

\begin{table}[ht]\centering\small
\begin{tabular}{@{}lll@{}}
\toprule
result & Lean name & remarks\\
\midrule
\Cref{prop:floorsum} & \code{eval\_eq\_floorSum} & \\
\Cref{thm:carry} & \code{eval\_add}, \code{floorSum\_add} & two summands$^{*}$\\
\Cref{prop:additive} & \code{additive\_iff\_layout} & the equivalence\\
\Cref{lem:expansion} & \code{fexp\_coeff}, \code{fexp\_unique} & existence, uniqueness$^{\dagger}$\\
\Cref{thm:spectrum} & \code{additive\_iff\_spec} & shapes without unit modes$^{\S}$\\
\Cref{thm:spectrum}(a) & \code{layout\_iff\_chain} & \\
\Cref{cor:coalesce}, merge step & \code{eval\_merge} & \\
\Cref{thm:defect} & \code{eval\_dot}, \code{eval\_comp} & coordinates; indices\\
\Cref{cor:naive} & \code{eval\_comp\_eq\_naive\_iff} & pointwise, both directions\\
\Cref{lem:forced}, $T'=T$ & \code{comp\_over\_shape\_unique} & \\
\Cref{lem:forced} with \cref{cor:naive} & \code{comp\_over\_shape\_iff} & see below$^{\ddagger}$\\
\Cref{thm:carryfree}(a) & \code{eval\_comp\_of\_carryFree\_ineq} & \\
\Cref{thm:carryfree}(b) & \code{carryFree\_of\_naive\_correct} & first assertion$^{\P}$\\
\Cref{ex:docs,ex:shah,ex:three} & file \code{Examples.lean} & by evaluation\\
\Cref{lem:swizzle}(a), (b) & \code{swizzle\_involutive}, \code{swizzle\_xor} & see below$^{\sharp}$\\
\bottomrule
\end{tabular}
\caption{Machine-checked statements. $^{*}$The case of $n$ summands is contained in
\code{eval\_dot}. $^{\dagger}$With the coefficients defined by the triangular recursion; the M\"obius
formula is not formalised. $^{\S}$More generally \code{layout\_over\_shape\_iff}, for shapes with pairwise distinct scales and
any initial segment of the domain; parts (b) and (c) are not formalised separately. $^{\ddagger}$A layout with the shape of $B$ and function $g$ exists if and only if the
carry defect vanishes on the whole domain. $^{\P}$The naive composite is correct on the whole
domain only if $B$ is carry-free; \code{naive\_correct\_iff\_carryFree} combines (a) and (b).
$^{\sharp}$Only the involution in (a) and the XOR identity in (b), for one field and a right
shift; the integer formula and the layout representation are not formalised.}\label{tab:lean}
\end{table}

Not formalised, and resting on the proofs given in the text together with the checks of
\cref{sec:checks}: the uniqueness of the strides in \cref{prop:additive}, the uniqueness statement
of \cref{cor:coalesce},
\cref{cor:recognition}, \cref{lem:forced} for proper refinements, \cref{thm:exact}, the last
assertion of \cref{thm:carryfree}(b), \cref{cor:universal},
\cref{prop:monomial,prop:dilated,prop:laws}, all of \cref{sec:complement,sec:beyond}, and the
rest of \cref{sec:swizzle}.

\subsection{Checks of the theory}\label{sec:checks}

The criteria and examples of \cref{sec:carry,sec:normalform,sec:composition,sec:complement} are
checked by a program that is independent of all CuTe implementations: \code{layout\_theory.py} implements
the definitions (evaluation, floor spectrum, recognition, exact composition, complement), and
\code{verify\_theory.py} runs fifteen checks, each exhaustive over a small universe or
randomised. For instance, the spectral criterion (\cref{thm:spectrum}) is compared with a direct
search over all shapes for all $14\,209$ functions on $\rng N$, $N\le9$, with values in a small
range; the composition formula and the four conditions of \cref{thm:exact} are compared with a
brute-force search over all by-mode refinements on $6\,000$ random pairs; \cref{prop:dilated} is
checked for all $50\,826$ triples $(q,\rho,t)$ with $q<40$; Long's theorem for all $M\le24$.
The same script tests the published sufficient conditions on random instances: the multi-mode
condition of~\cite{shah2024} was sound in all $26\,798$ admissible instances with at most two
modes and failed in $21$ of $9\,124$ admissible instances with three (\cref{ex:three} is one), and
the segregation condition of~\cite{cecka2026}, with stride divisibility taken with respect to
all scales of the outer layout, excluded interference between modes in all $26\,436$ instances
that satisfy it.
A draft of \cref{sec:background,sec:carry,sec:normalform,sec:composition,sec:complement,sec:beyond}
was in addition reviewed independently, with separately written exhaustive checks; the errors
found are corrected in this version.
The statements of \cref{sec:swizzle} are checked by a second script, \code{verify\_swizzle.py}
(fourteen checks, for CuTe's swizzles and for arbitrary pairings): for instance \cref{cor:xorlinear} against a search over all by-mode refinements
for $6\,000$ random layouts, \cref{thm:decay} on $20\,000$ random bit-disjoint layouts with
offsets, \cref{prop:fields} on $30\,000$ random instances (with counterexamples when a hypothesis
is dropped), and the merging rule on all $360$ admissible pairs of swizzles with $1\le B\le3$,
$M\le3$, $|S|\le4$.

\subsection{Differential testing of implementations}\label{sec:testing}

We compared the theory with three implementations of \code{composition} and the operations built
on it:
\begin{itemize}[leftmargin=1.4em,itemsep=1pt]
\item \textbf{C++ CuTe}, the headers of CUTLASS at commit \code{af95104} (20 November 2026), with
 static (compile-time) integers, compiled for the host with \code{clang} against NVIDIA's
 redistributable CUDA~13.5 headers. Inadmissible compositions are compile errors here.
\item \textbf{pycute}, the Python module shipped in the same CUTLASS checkout.
\item \textbf{PyCuTe}~\cite{pycute}, the reference implementation accompanying~\cite{cecka2026}, at
 commit \code{56cc3d0}.
\end{itemize}
Test inputs are small flat layouts drawn at random (outer layout: rank $\le3$, mode sizes
$\le 8$, strides in $\{0,\dots,24\}$; inner layout: rank $\le 3$, sizes $\le6$, strides $\le12$),
the same layouts with random nesting, and a ``structured'' family in which the outer layout is a
column-major, row-major, padded or digit-permuted layout of some size $M\le64$ and the inner one
is a compact layout of size $M$ or a tile obtained from one by dropping modes. An answer counts as
correct if it is a layout of the right size whose values equal $\ext A(B(i))$ pointwise; a
rejection is correct if no composite exists in the sense of \cref{def:composition}. The
distributions are arbitrary: the counts in \cref{tab:difftest} show that the phenomena exist and
how they split by cause, not how often they occur in practice.

\begin{table}[t]\centering\small
\begin{tabular}{@{}llrrrrr@{}}
\toprule
implementation & inputs & cases & correct & \shortstack{correctly\\rejected} & \shortstack{silently\\wrong} & \shortstack{rejected,\\composable}\\
\midrule
C++ CuTe (static) & random & 3\,000 & 1\,846 & 623 & 153 & 378\\
C++ CuTe (static) & structured & 3\,000 & 2\,579 & 337 & 50 & 34\\
PyCuTe & random & 20\,000 & 12\,925 & 4\,665 & 352 & 2\,058\\
PyCuTe & random, nested & 20\,000 & 12\,930 & 4\,728 & 323 & 2\,019\\
PyCuTe & structured & 20\,000 & 17\,158 & 2\,624 & 36 & 182\\
pycute & random & 20\,000 & 10\,505 & 4\,153 & 2\,072 & 3\,270\\
pycute & structured & 20\,000 & 17\,121 & 2\,466 & 194 & 219\\
\bottomrule
\end{tabular}
\caption{Outcomes of \code{composition(A,B)} against pointwise evaluation.}\label{tab:difftest}
\end{table}

\paragraph{Silently wrong results.}
In PyCuTe every wrong result ($352+323+36$) is an instance of interference between modes
(\cref{sec:deviations}): each mode of the inner layout is composed correctly, the concatenation
is not the composite, and in fact no composite exists, so the call should have been rejected.
The same happens in C++. There, in addition, the static stride check accepts a stride that is
smaller than a mode size without requiring that it divide it, so some \emph{single-mode}
compositions are wrong as well: of the $153$ wrong results on random inputs, $53$ are pure
interference, $29$ have a single-mode inner layout, and $71$ contain a wrong single mode. The
module pycute further lacks the shape-divisibility check, and returns layouts of the wrong size.
\Cref{tab:examples} lists small instances; all were confirmed by compiling and running them.

\begin{table}[t]\centering\small
\begin{tabular}{@{}llll@{}}
\toprule
$A$ & $B$ & C++ result & $\ext A(B(i))$\\
\midrule
$(4,3):(3,1)$ & $(3,4):(1,3)$ & $(3,(2,2)):(3,(9,1))$ & $0,3,6,9,1,4,7,10,2,5,8,11$\\
$(5,2):(1,8)$ & $4:4$ & $(2,2):(4,8)$ & $0,4,11,18$\\
$(3,4):(1,5)$ & $(2,2):(1,2)$ & $(2,2):(1,2)$ & $0,1,2,5$\\
$(6,2):(1,7)$ & $(3,2):(2,3)$ & $(3,2):(2,3)$ & $0,2,4,3,5,8$\\
$(12,2):(1,13)$ & $(2,2,2):(3,4,6)$ & $(2,2,2):(3,4,6)$ & $0,3,4,7,6,9,10,14$\\
\bottomrule
\end{tabular}
\caption{Compositions that compile with static integers in C++ CuTe and return a layout whose
values differ from $\ext A\circ\bnd B$. The first views a row-major $4\times3$ matrix through the
column-major shape $(3,4)$; the third views the first four elements of a padded $3\times4$ matrix
as $2\times2$. PyCuTe rejects the first two and returns the same wrong layouts for the last
three. In none of the five cases does a composite exist.}\label{tab:examples}
\end{table}

\paragraph{Rejected although composable.}
Typical causes are those predicted in \cref{sec:deviations}: strides with several non-zero digits
(``sum strides''), and single digits that do not divide the mode size but stay within the range
where \cref{prop:dilated} applies. Most instances have a very small inner layout (for example
$(2,1):(3,12)\circ 2:3=2:15$), and we do not claim that they matter in practice; they show that
the divisibility conditions are sufficient and not necessary.

\paragraph{Other operations (PyCuTe).}
\code{coalesce} returned the normal form of \cref{cor:coalesce} in all $20\,000$ cases (up to
the stride of a layout of size $1$). \code{complement} agreed with \cref{thm:complement} in all
$4\,382$ cases in which a complement in $\rng M$ exists, rounded $M$ up in $7\,239$ cases and
returned a non-tiling layout in $1\,969$ (\cref{rem:nontiling}). \code{logical\_divide} was correct
in $14\,346$ cases and correctly rejected $5\,617$; in $33$ further cases its result had a
different number of elements than $A$, and in $4$ it rejected a divisible pair. These appear to
share one cause: the complement is extended along the shape of $A$ as given, not of $\coal(A)$, so
that, for instance, \code{logical\_divide} of $(3,4):(4,12)$---the layout $12:4$ written with two
modes---by $4:1$ returns $(4,2):(4,16)$, covering $8$ of the $12$ elements (C++ CuTe coalesces
first and returns $(4,3):(4,16)$). \code{logical\_product}
was correct in $8\,915$ cases and silently wrong in $25$, again by interference, inside the
composition $A^*\circ B$. The inverses satisfied their post-conditions in all cases.

\paragraph{Swizzles.}
The operations of \cref{sec:swzops} were tested against C++ CuTe on random instances with small
parameters ($B\le2$, $M\le3$; inner layouts with at most three modes for slices and at most five
binary modes for \code{composition(a, Swizzle)}; initial offset $0$; coordinates of slices static
or run-time). \Cref{tab:swz} gives the outcomes and
\cref{tab:swzex} small instances. For \code{slice\_and\_offset}, $286$ of the $600$ inner layouts
are bit-disjoint in the sense of \cref{def:bitdisjoint}; for these, $270$ slices were correct and
$16$ did not compile. All $29$ wrong results came from inner layouts that are not
bit-disjoint---modes that share a bit, or a stride such as $6$ whose multiples overlap in
binary---as \cref{prop:fields,thm:decay} predict; in $5$ cases CuTe kept a swizzled layout
although the slice decays. For \code{composition(a, Swizzle)}, in none of the $13$ wrong cases is
there any \code{Swizzle<B,M,S>} that represents $a\circ\sigma$, so these calls should have been
rejected; conversely $19$ of the rejected cases are representable. The merging rule behaved as
stated in \cref{sec:swzops} in all $57$ cases, and \code{upcast} was correct in all $600$ (all
offsets there are multiples of the cast factor). \code{cosize} differed from the largest offset
plus one in $198$ of $600$ cases and was too small in $149$ of them; for compact inner layouts of
power-of-two size it was exact in $294$ cases and too small in $20$, which by
\cref{lem:swizzle}(d) can only happen for $S<0$.

In PyCuTe, \code{composition} of an $\F_2$-stride layout $A$ with power-of-two mode sizes and an
integer layout $B$ was tested on $20\,000$ pairs. For the $14\,165$ pairs with a bit-disjoint $B$
it returned the correct layout in $13\,828$ and rejected the other $337$; none was wrong. For the
remaining $5\,835$ it was correct in $645$, correctly rejected $2\,256$, rejected $31$ composable
pairs, and returned a wrong layout in $2\,903$: in $2\,301$ of these every mode is composed
correctly and the modes interfere, in $602$ a single mode is already wrong.
\code{right\_inverse} satisfied its post-condition on all $4\,000$ $\F_2$-stride layouts tested.
CUTLASS's \code{pycute} only evaluates swizzled layouts.

\begin{table}[t]\centering\small
\begin{tabular}{@{}lrrrr@{}}
\toprule
operation (C++ CuTe) & cases & rejected & correct & wrong\\
\midrule
\code{slice\_and\_offset} of a swizzled layout & 600 & 207 & 364 & 29\\
\code{composition(a, Swizzle)} & 600 & 201 & 386 & 13\\
\code{right\_inverse} of a swizzled layout & 600 & 66 & 530 & 4\\
\code{composition(Swizzle, Swizzle)} & 57 & 22 & 35 & 0\\
\code{upcast} of a swizzled layout & 600 & 0 & 600 & 0\\
\code{cosize} of a swizzled layout & 600 & 0 & 402 & 198\\
\bottomrule
\end{tabular}
\caption{Operations on swizzled layouts in C++ CuTe against pointwise evaluation. ``Rejected''
means that the call does not compile.}\label{tab:swz}
\end{table}

\begin{table}[t]\centering\small
\begin{tabular}{@{}p{0.36\linewidth}p{0.33\linewidth}p{0.24\linewidth}@{}}
\toprule
call & C++ result and its values & correct values\\
\midrule
slice $(1,\_,\_)$ of \code{Sw<1,1,1>}${}\circ(2,2,2):(1,1,4)$ & $(2,2):(1,6)$ plus $1$:\newline $1,2,7,8$ & $1,2,7,4$\\[2pt]
slice $(1,\_)$ of \code{Sw<2,1,2>}${}\circ(2,4):(1,6)$ & $(2,2):(6,10)$ plus $1$:\newline $1,7,11,17$ & $1,7,15,23$\\[2pt]
slice $(1,\_)$ of \code{Sw<1,0,1>}${}\circ(2,2):(1,2)$\newline with offset $1$ & $2:3$ plus $0$:\newline $0,3$ & $3,4$\\[2pt]
$(2,2,4):(8,4,1)\circ{}$\code{Sw<2,0,2>} & \code{Sw<2,0,-2>}${}\circ(2,2,4):(8,4,1)$:\newline $0,8,4,12,5,13,1,9,\dots$ & $0,8,4,12,9,1,13,5,\dots$\\[2pt]
\code{right\_inverse}$(L)$ for\newline $L={}$\code{Sw<2,0,2>}${}\circ(4,2,2):(1,8,4)$ & $R={}$\code{Sw<2,0,2>}${}\circ(4,2,2):(1,8,4)$:\newline $L(R(k))=0,1,2,3,7,6,5,4,\dots$ & $L(R(k))=k$\\
\bottomrule
\end{tabular}
\caption{Operations on swizzled layouts that compile in C++ CuTe and return a wrong result
(\code{Sw<B,M,S>} abbreviates \code{Swizzle<B,M,S>}; the correct values are $\sigma(o+f(c))$
for the slices and $a(\sigma(i))$ in the fourth row). In the first two rows the inner layout is
not bit-disjoint; in the third, which is not one of the random instances, the initial offset
overlaps the fixed coordinate; in the last two
the layout maps the two fields of the swizzle with different bit orders.}\label{tab:swzex}
\end{table}

\paragraph{Threats to validity.}
The layouts are small and the distributions synthetic; C++ was exercised with static layouts
only (with dynamic integers CuTe performs no admissibility checks, by design) and on a host
compiler; swizzle fields were confined to low bits and swizzled layouts start from offset $0$; all three implementations were tested at one commit each.
